\documentclass[UTF8,12pt]{ctexart}
\usepackage[a4paper,margin=24mm]{geometry}
\usepackage{amsmath,amssymb,bm,graphicx,adjustbox}
\newcommand{\fitmath}[1]{\begin{adjustbox}{max width=\linewidth}$\displaystyle #1$\end{adjustbox}}
\setlength{\parindent}{0pt}
\setlength{\parskip}{.7em}
\sloppy
\allowdisplaybreaks
\begin{document}
\section*{1997 年考研数学三 · 参考答案}
\subsection*{1997 年真题参考答案}

\subsection*{一、填空题}

（1）\(\displaystyle e^{f(x)}\left[\dfrac{\displaystyle 1}{\displaystyle x} f^{\prime}(\ln x)+f^{\prime}(x)f(\ln x)\right]\,dx\)。 （2）\(\displaystyle \dfrac{\displaystyle \pi}{\displaystyle 4-\pi}\)。 （3）\(\displaystyle y_t=C+(t-2)2^t\)，其中 \(\displaystyle C\) 为任意常数。 （4）\(\displaystyle -\sqrt2<t<\sqrt2\)。 （5）\(\displaystyle t\) 分布；参数为 \(\displaystyle 9\)。


\subsection*{二、选择题}

（1）B。 （2）C。 （3）C。 （4）D。 （5）A。


三、证明略。


四、\(\displaystyle \dfrac{\displaystyle du}{\displaystyle dx}=\dfrac{\displaystyle \partial f}{\displaystyle \partial x}+\dfrac{\displaystyle y^2}{\displaystyle 1-xy}\dfrac{\displaystyle \partial f}{\displaystyle \partial y}+\dfrac{\displaystyle z}{\displaystyle xz-x}\dfrac{\displaystyle \partial f}{\displaystyle \partial z}\)。


五、（1）当销售量 \(\displaystyle x=\dfrac{\displaystyle 5}{\displaystyle 2}(4-t)\) 时，利润最大。（2）\(\displaystyle t=2\) 时，政府税收总额最大。


六、证明略。（证明 \(\displaystyle \lim_{x\to0^+}F(x)=F(0)\)，并且当 \(\displaystyle x>0\) 时，\(\displaystyle F^{\prime}(x)\geq0\)。）


七、（1）\(\displaystyle \overline{OP_n}=\dfrac{\displaystyle 1}{\displaystyle 2^{n-1}}\)。（2）\(\displaystyle \dfrac{\displaystyle 4}{\displaystyle 3}\)。


八、\(\displaystyle f(t)=(4\pi t^2+1)e^{4\pi t^2}\)。


九、（1）\(\displaystyle PQ=\begin{pmatrix}A&\alpha\\0&|A|(b-\alpha^{\mathrm T}A^{-1}\alpha)\end{pmatrix}\)，其中零向量 \(\displaystyle 0\) 为 \(\displaystyle n\) 维行向量。（2）证明略。（利用（1）的结论，证明 \(\displaystyle |Q|\ne0\iff b-\alpha^{\mathrm T}A^{-1}\alpha\ne0\)。）


十、（1）\(\displaystyle A\) 的属于特征值 3 的特征向量为 \(\displaystyle k(1,\allowbreak 0,\allowbreak 1)^{\mathrm T}\)，其中 \(\displaystyle k\) 为任意非零常数。（2）\(\displaystyle A=\dfrac{\displaystyle 1}{\displaystyle 6}\begin{pmatrix}13&-2&5\\-2&10&2\\5&2&13\end{pmatrix}\)。


十一、\(\displaystyle F(x)=\begin{cases}0,\allowbreak &x<-1,\allowbreak \\\dfrac{\displaystyle 5x+7}{\displaystyle 16},\allowbreak &-1\leq x<1,\allowbreak \\1,\allowbreak &x\geq1.\end{cases}\)


十二、\(\displaystyle \dfrac{\displaystyle 70}{\displaystyle 6}\approx11.67\)。


十三、\(\displaystyle f(t)=\begin{cases}25te^{-5t},\allowbreak &t>0,\allowbreak \\0,\allowbreak &t\leq0.\end{cases}\) 数学期望为 \(\displaystyle \dfrac{\displaystyle 2}{\displaystyle 5}\)，方差为 \(\displaystyle \dfrac{\displaystyle 2}{\displaystyle 25}\)。

\end{document}
